\documentclass[aspectratio=169,11pt]{beamer}

% Compatibilidad con compilaciones mediante pdfLaTeX.
\usepackage[utf8]{inputenc}
\usepackage[T1]{fontenc}
\usepackage[spanish]{babel}

% Matemáticas y diagramas.
\usepackage{amsmath,amssymb,mathtools}
\usepackage{tikz}
\usetikzlibrary{calc,positioning}
\usepackage{tikz-cd}

% Tema docente, incorporado como submódulo en beamerteaching/.
\makeatletter
\def\input@path{{beamerteaching/}}
\makeatother
\usetheme{Teaching}
\titlegraphic{%
  \includegraphics[width=0.20\paperwidth]{beamerteaching/assets/logo-us.png}%
}

\title{El teorema de escisión}
\subtitle{Homología singular}
\author{Fernando Muro}
\institute{Universidad de Sevilla}
\date{}

% Comandos generales.
\newcommand{\N}{\mathbb{N}}
\newcommand{\Z}{\mathbb{Z}}
\newcommand{\R}{\mathbb{R}}
\newcommand{\id}{\operatorname{id}}
\newcommand{\im}{\operatorname{im}}
\newcommand{\diam}{\operatorname{diam}}
\newcommand{\Int}{\operatorname{int}}           % interior de un subconjunto

% Homología singular.
\newcommand{\simplex}[1]{\Delta^{#1}}          % símplice estándar
\newcommand{\Sing}{C}                           % cadenas singulares C_n(X)
\newcommand{\LC}{\mathit{LC}}                   % cadenas singulares lineales LC_n(Y)
\newcommand{\cU}{\mathcal{U}}                   % recubrimiento {U_i}_{i\in I}
\newcommand{\CU}[1][\cU]{\Sing^{\mkern2mu#1}}              % cadenas U-pequeñas C^U_n(X)
\newcommand{\tCU}[1][\cU]{\widetilde{\Sing}^{\mkern2mu#1}}      % versión aumentada
\newcommand{\Sd}{S}                             % operador subdivisión en cadenas
\newcommand{\sd}{\operatorname{sd}}             % familia de símplices de la subdivisión

% Colores predefinidos.
\usebeamercolor{alerted text}
\colorlet{alertcolor}{fg}
\usebeamercolor{frame rule}
\colorlet{secondarycolor}{bg}

\begin{document}

\begin{frame}
  \titlepage
\end{frame}

\section{Subdivisión baricéntrica de un símplice}

% ------------------------------------------------------------------------
\subsection{Símplices generalizados}

\begin{frame}
  \frametitle{Símplices generalizados}

  \begin{definition}
    Dada una lista de $n+1$ puntos $v_0,\ldots,v_n\in\R^N$, el
    \alert{símplice generalizado} de dimensión $n$ y vértices
    $v_0,\ldots,v_n$ es su envolvente convexa
    \[
      [v_0,\ldots,v_n]=\Bigl\{\,t_0v_0+\cdots+t_nv_n \;:\;
      t_i\geq0,\ t_0+\cdots+t_n=1\,\Bigr\}.
    \]
    Los $(t_0,\ldots,t_n)$ son \alert{unas coordenadas baricéntricas} del punto
    $t_0v_0+\cdots+t_nv_n$.
  \end{definition}

  \medskip

  Los vértices pueden repetirse o puede haber tres alineados.
  El símplice generalizado viene dado por la
  \alert{lista} de sus vértices, en ese orden: listas distintas dan
  símplices generalizados distintos, aunque tengan la misma envolvente.
\end{frame}

\begin{frame}
  \frametitle{Símplices y símplices degenerados}

  Hay dos tipos de símplices generalizados $[v_0,\ldots,v_n]$:
  \begin{itemize}
    \item si $v_0,\ldots,v_n$ son afínmente independientes, es un
      \alert{símplice} en el sentido usual, y cada punto tiene
      coordenadas baricéntricas \alert{únicas};
    \item en otro caso, decimos que es un \alert{símplice degenerado}.
  \end{itemize}

  \medskip

  En lo que sigue, «símplice» a secas significa siempre no degenerado.
\end{frame}

\begin{frame}
  \frametitle{Ejemplos de símplices}

  \begin{columns}[T]
    \begin{column}{0.30\textwidth}
      \centering
      \begin{tikzpicture}
        \useasboundingbox (-0.5,-0.9) rectangle (2.9,2.4);
        \coordinate (A) at (0,0.6);
        \coordinate (B) at (2.4,0.6);
        \draw[thick] (A) -- (B);
        \fill (A) circle (1.6pt) node[below] {$v_0$};
        \fill (B) circle (1.6pt) node[below] {$v_1$};
      \end{tikzpicture}

      \small Un segmento $[v_0,v_1]$.
    \end{column}
    \begin{column}{0.30\textwidth}
      \centering
      \begin{tikzpicture}
        \useasboundingbox (-0.5,-0.9) rectangle (2.9,2.4);
        \coordinate (A) at (0,0);
        \coordinate (B) at (2.4,0);
        \coordinate (C) at (1.2,2);
        \filldraw[thick,fill=alertcolor!10] (A) -- (B) -- (C) -- cycle;
        \fill (A) circle (1.6pt) node[below] {$v_0$};
        \fill (B) circle (1.6pt) node[below] {$v_1$};
        \fill (C) circle (1.6pt) node[above] {$v_2$};
      \end{tikzpicture}

      \small Un triángulo $[v_0,v_1,v_2]$.
    \end{column}
    \begin{column}{0.30\textwidth}
      \centering
      \begin{tikzpicture}
        \useasboundingbox (-0.5,-0.9) rectangle (2.9,2.4);
        \coordinate (A) at (0,0.1);
        \coordinate (B) at (1.7,-0.4);
        \coordinate (C) at (2.5,0.6);
        \coordinate (D) at (1.2,2);
        \fill[alertcolor!10] (A) -- (B) -- (C) -- (D) -- cycle;
        \draw[dashed] (A) -- (C);
        \draw[thick] (A) -- (B) -- (C) -- (D) -- cycle (B) -- (D);
        \fill (A) circle (1.6pt) node[left] {$v_0$};
        \fill (B) circle (1.6pt) node[below] {$v_1$};
        \fill (C) circle (1.6pt) node[right] {$v_2$};
        \fill (D) circle (1.6pt) node[above] {$v_3$};
      \end{tikzpicture}

      \small Un tetraedro $[v_0,v_1,v_2,v_3]$.
    \end{column}
  \end{columns}

  \bigskip

  En todos ellos, coordenadas baricéntricas distintas dan puntos distintos:
  cada punto tiene coordenadas baricéntricas \alert{únicas}.
\end{frame}

\begin{frame}
  \frametitle{Ejemplos de símplices degenerados}

  \begin{columns}[T]
    \begin{column}{0.22\textwidth}
      \centering
      \begin{tikzpicture}
        \useasboundingbox (-1.1,-0.9) rectangle (1.1,2.2);
        \fill (0,0.6) circle (1.8pt) node[below=4pt] {$v_0=v_1$};
      \end{tikzpicture}

      \footnotesize
      Si $v_0=v_1$, $[v_0,v_1]$ es un punto. Todo $(t_0,t_1)$ da el mismo
      punto.
    \end{column}
    \begin{column}{0.24\textwidth}
      \centering
      \begin{tikzpicture}
        \useasboundingbox (-0.6,-0.9) rectangle (2.4,2.2);
        \coordinate (A) at (0,0.6);
        \coordinate (B) at (2,0.6);
        \draw[thick] (A) -- (B);
        \fill (A) circle (1.6pt) node[below=4pt] {$v_0=v_1$};
        \fill (B) circle (1.6pt) node[below=4pt] {$v_2$};
      \end{tikzpicture}

      \footnotesize
      Si $v_0=v_1$, $[v_0,v_1,v_2]$ es un segmento. El punto $v_0$ tiene
      coordenadas $(1,0,0)$ y $(0,1,0)$.
    \end{column}
    \begin{column}{0.24\textwidth}
      \centering
      \begin{tikzpicture}
        \useasboundingbox (-0.4,-0.9) rectangle (2.4,2.2);
        \coordinate (A) at (0,0.6);
        \coordinate (B) at (2,0.6);
        \coordinate (M) at (1,0.6);
        \draw[thick] (A) -- (B);
        \fill (A) circle (1.6pt) node[below=4pt] {$v_0$};
        \fill (B) circle (1.6pt) node[below=4pt] {$v_1$};
        \fill (M) circle (1.6pt) node[below=4pt] {$v_2$};
      \end{tikzpicture}

      \footnotesize
      Si $v_2=\frac12v_0+\frac12v_1$, $[v_0,v_1,v_2]$ es un segmento.
      El punto $v_2$ tiene coordenadas $(0,0,1)$ y $(\frac12,\frac12,0)$.
    \end{column}
    \begin{column}{0.24\textwidth}
      \centering
      \begin{tikzpicture}
        \useasboundingbox (-0.5,-0.9) rectangle (2.5,2.2);
        \coordinate (A) at (0,0);
        \coordinate (B) at (2,0);
        \coordinate (C) at (2.2,1.6);
        \coordinate (D) at (0.3,1.4);
        \filldraw[thick,fill=alertcolor!10] (A) -- (B) -- (C) -- (D) -- cycle;
        \fill (A) circle (1.6pt) node[below] {$v_0$};
        \fill (B) circle (1.6pt) node[below] {$v_1$};
        \fill (C) circle (1.6pt) node[above] {$v_2$};
        \fill (D) circle (1.6pt) node[above] {$v_3$};
      \end{tikzpicture}

      \footnotesize
      Con vértices coplanarios, $[v_0,v_1,v_2,v_3]$ es un cuadrilátero, de dimensión $3$
      como símplice generalizado.
    \end{column}
  \end{columns}
\end{frame}

\begin{frame}
  \frametitle{Caras y diámetro}

  \begin{definition}
    Sea $\sigma=[v_0,\ldots,v_n]$ un símplice generalizado con $n\geq1$. Su
    \alert{cara $i$-ésima} es el símplice generalizado de dimensión $n-1$
    \[
      \partial_i\sigma=[v_0,\ldots,\widehat{v_i},\ldots,v_n],
      \qquad 0\leq i\leq n,
    \]
    obtenido al suprimir el vértice $v_i$.

    \medskip

    El \alert{diámetro} de $\sigma$ es
    \[
      \diam\sigma=\sup\{\,|x-y| : x,y\in\sigma\,\}.
    \]
  \end{definition}

  \medskip

  Cada cara está contenida en $\sigma$, luego
  $\diam \partial_i\sigma\leq\diam\sigma$. Las caras de un símplice son símplices;
  las de un símplice degenerado pueden serlo o no.
\end{frame}

% ------------------------------------------------------------------------
\subsection{Baricentro y subdivisión}

\begin{frame}
  \frametitle{El baricentro}

  \begin{definition}
    El \alert{baricentro} del símplice generalizado $\sigma=[v_0,\ldots,v_n]$
    es el punto de coordenadas baricéntricas
    $\bigl(\frac{1}{n+1},\ldots,\frac{1}{n+1}\bigr)$:
    \[
      b_\sigma=\frac{1}{n+1}\,(v_0+v_1+\cdots+v_n)\in\sigma.
    \]
  \end{definition}

  \begin{columns}[c]
    \begin{column}{0.58\textwidth}
      Aunque se pierda la unicidad de las coordenadas, el baricentro
      tiene sentido. Depende de la \alert{lista} de vértices, no solo de la
      envolvente: $[a,b]$ y $[a,a,b]$ son el mismo segmento, pero
      \[
        b_{[a,b]}=\tfrac12a+\tfrac12b,
        \qquad
        b_{[a,a,b]}=\tfrac23a+\tfrac13b.
      \]
    \end{column}
    \begin{column}{0.38\textwidth}
      \centering
      \begin{tikzpicture}
        \coordinate (A) at (0,0);
        \coordinate (B) at (3.6,0);
        \draw[thick] (A) -- (B);
        \fill (A) circle (1.6pt) node[below] {$a$};
        \fill (B) circle (1.6pt) node[below] {$b$};
        \fill[alertcolor] (1.8,0) circle (1.6pt)
          node[above] {$b_{[a,b]}$};
        \fill[alertcolor] (1.2,0) circle (1.6pt)
          node[below] {$b_{[a,a,b]}$};
      \end{tikzpicture}
    \end{column}
  \end{columns}
\end{frame}

\begin{frame}
  \frametitle{La subdivisión baricéntrica}

  \begin{definition}
    La \alert{subdivisión baricéntrica} de un símplice generalizado
    $\sigma$ de dimensión $n$ es la familia finita $\sd(\sigma)$ de
    símplices generalizados de dimensión $n$ definida por recurrencia en
    $n$:
    \begin{itemize}
      \item si $n=0$, $\sd(\sigma)=\{\sigma\}$;
      \item si $n\geq1$, $\sd(\sigma)$ está formada por los
        \[
          [b_\sigma,w_0,\ldots,w_{n-1}]
          \quad\text{con}\quad
          [w_0,\ldots,w_{n-1}]\in\sd(\partial_i\sigma),\ 0\leq i\leq n.
        \]
    \end{itemize}
  \end{definition}

  \medskip

  Es decir, se subdividen las caras y se hace el \alert{cono} desde el
  baricentro sobre cada símplice generalizado obtenido. En la familia puede
  haber repeticiones si $\sigma$ es degenerado.
\end{frame}

\begin{frame}
  \frametitle{Ejemplos}

  \begin{columns}[T]
    \begin{column}{0.36\textwidth}
      \centering
      \begin{tikzpicture}[scale=0.9]
        \coordinate (A) at (0,0);
        \coordinate (B) at (3,0);
        \coordinate (M) at (1.5,0);
        \draw[thick] (A) -- (B);
        \foreach \p/\l/\pos in {A/v_0/below, B/v_1/below, M/b_\sigma/above}
          \fill (\p) circle (1.6pt) node[\pos] {$\l$};
        \node at (1.5,-1.2) {$n=1$};
      \end{tikzpicture}

      \medskip

      \small
      Si $n=1$, $\sd[v_0,v_1]=\{[b_\sigma,v_1],[b_\sigma,v_0]\}$.
    \end{column}
    \begin{column}{0.60\textwidth}
      \centering
      \begin{tikzpicture}[scale=0.9]
        \coordinate (A) at (0,0);
        \coordinate (B) at (4,0);
        \coordinate (C) at (2,3.3);
        \coordinate (G) at ($1/3*(A)+1/3*(B)+1/3*(C)$);
        \draw (G) -- (A) (G) -- (B) (G) -- (C);
        \draw (G) -- ($(A)!0.5!(B)$) (G) -- ($(B)!0.5!(C)$)
              (G) -- ($(C)!0.5!(A)$);
        \draw[thick] (A) -- (B) -- (C) -- cycle;
        \foreach \p in {A,B,C,G}
          \fill (\p) circle (1.6pt);
        \foreach \p/\q in {A/B,B/C,C/A}
          \fill ($(\p)!0.5!(\q)$) circle (1.3pt);
        \node[below left] at (A) {$v_0$};
        \node[below right] at (B) {$v_1$};
        \node[above] at (C) {$v_2$};
        \node[left=5pt] at (G) {$b_\sigma$};
      \end{tikzpicture}

      \small
      Si $n=2$, hay seis triángulos $[b_\sigma,b_{\partial_i\sigma},v_j]$ con $j\neq i$.
    \end{column}
  \end{columns}

  \medskip

  En general $\sd(\sigma)$ tiene $(n+1)!$ elementos, y si $\sigma$ es un
  símplice, estos son símplices y recubren $\sigma$. No usaremos estos
  hechos.
\end{frame}

% Dibuja la subdivisión baricéntrica del triángulo de vértices #1, #2, #3.
% Se define fuera del frame: Beamer no admite # en el cuerpo de un frame.
\newcommand{\subdividir}[3]{%
  \coordinate (g) at ($1/3*(#1)+1/3*(#2)+1/3*(#3)$);
  \draw[very thin] (g) -- (#1) (g) -- (#2) (g) -- (#3)
    (g) -- ($(#1)!0.5!(#2)$) (g) -- ($(#2)!0.5!(#3)$)
    (g) -- ($(#3)!0.5!(#1)$);}

\begin{frame}
  \frametitle{Subdivisión iterada}

  \begin{definition}
    Para $k\geq0$ definimos $\sd^0(\sigma)=\{\sigma\}$ y
    \[
      \sd^{k+1}(\sigma)=\bigcup_{\tau\in\sd^k(\sigma)}\sd(\tau).
    \]
  \end{definition}

  \begin{center}
    \begin{tikzpicture}[scale=0.85]
      \coordinate (A) at (0,0);
      \coordinate (B) at (4,0);
      \coordinate (C) at (2,3.3);
      \coordinate (G) at ($1/3*(A)+1/3*(B)+1/3*(C)$);
      \coordinate (AB) at ($(A)!0.5!(B)$);
      \coordinate (BC) at ($(B)!0.5!(C)$);
      \coordinate (CA) at ($(C)!0.5!(A)$);
      \subdividir{G}{A}{AB}
      \subdividir{G}{AB}{B}
      \subdividir{G}{B}{BC}
      \subdividir{G}{BC}{C}
      \subdividir{G}{C}{CA}
      \subdividir{G}{CA}{A}
      \draw (G) -- (A) (G) -- (B) (G) -- (C) (G) -- (AB) (G) -- (BC)
        (G) -- (CA);
      \draw[thick] (A) -- (B) -- (C) -- cycle;
    \end{tikzpicture}
  \end{center}

  \centering
  $\sd^2(\sigma)$ para un triángulo: los símplices se hacen
  \alert{cada vez más pequeños}.
\end{frame}

% ------------------------------------------------------------------------
\subsection{El diámetro de la subdivisión}

\begin{frame}
  \frametitle{La subdivisión está dentro de $\sigma$}

  \begin{lemma}
    Si $\tau\in\sd(\sigma)$, entonces $\tau\subseteq\sigma$.
  \end{lemma}

  \begin{proof}
    Como $\tau$ es la envolvente convexa de sus vértices, basta ver que los
    vértices de $\tau$ están en el convexo $\sigma$. Razonamos por
    inducción en $n$; el caso $n=0$ es trivial.

    \medskip

    Si $n\geq1$, entonces $\tau=[b_\sigma,w_0,\ldots,w_{n-1}]$ con
    $[w_0,\ldots,w_{n-1}]\in\sd(\partial_i\sigma)$. Ya sabemos que
    $b_\sigma\in\sigma$. Por hipótesis de inducción, cada $w_j$ está en
    $\partial_i\sigma$, que está contenida en $\sigma$.
  \end{proof}

  Por inducción en $k$, lo mismo vale para todo $\tau\in\sd^k(\sigma)$.
\end{frame}

\begin{frame}
  \frametitle{El diámetro lo determinan los vértices}

  \begin{lemma}
    Sea $\sigma=[v_0,\ldots,v_n]$. Para todo $p\in\R^N$ y todo
    $x\in\sigma$,
    \[
      |x-p|\leq\max_j|v_j-p|.
    \]
    En consecuencia, $\diam\sigma=\max_{i,j}|v_i-v_j|$.
  \end{lemma}

  \begin{proof}
    Si $x=\sum_jt_jv_j$ con $t_j\geq0$ y $\sum_jt_j=1$, entonces
    \[
      |x-p|=\Bigl|\sum_jt_j(v_j-p)\Bigr|
      \leq\sum_jt_j|v_j-p|
      \leq\max_j|v_j-p|.
    \]
    Si $x,y\in\sigma$, aplicando esto con $p=y$ y luego con $p=v_j$ se
    obtiene $|x-y|\leq\max_j|v_j-y|\leq\max_{i,j}|v_i-v_j|$. La
    desigualdad contraria se da porque los vértices están en $\sigma$.
  \end{proof}
\end{frame}

\begin{frame}
  \frametitle{Distancia al baricentro}

  \begin{lemma}
    Sea $\sigma=[v_0,\ldots,v_n]$. Para todo $x\in\sigma$,
    \[
      |x-b_\sigma|\leq\frac{n}{n+1}\diam\sigma.
    \]
  \end{lemma}

  \begin{proof}
    Por el lema anterior con $p=b_\sigma$, basta probarlo para los vértices $x=v_i$.
    Como
    \[
      v_i-b_\sigma=\frac{1}{n+1}\sum_{j=0}^n(v_i-v_j)
    \]
    y el sumando $j=i$ es nulo, quedan a lo sumo $n$ sumandos, cada uno de
    norma $\leq\diam\sigma$. Por tanto,
    $|v_i-b_\sigma|\leq\frac{n}{n+1}\diam\sigma$.
  \end{proof}
\end{frame}

\begin{frame}[allowframebreaks]
  \frametitle{El diámetro decrece}

  \begin{proposition}
    Sea $\sigma$ un símplice generalizado de dimensión $n$. Para todo
    $\tau\in\sd(\sigma)$,
    \[
      \diam\tau\leq\frac{n}{n+1}\diam\sigma.
    \]
  \end{proposition}

  La constante $\frac{n}{n+1}<1$ depende \alert{solo de la dimensión}.

  \begin{proof}
    Por inducción en $n$. Si $n=0$, ambos diámetros son nulos.

    \medskip

    Sea $n\geq1$ y $\tau=[b_\sigma,w_0,\ldots,w_{n-1}]$ con
    $[w_0,\ldots,w_{n-1}]\in\sd(\partial_i\sigma)$. Por el lema de los vértices,
    basta acotar las distancias entre vértices de $\tau$.

    \framebreak

    \begin{itemize}
      \item \alert{Entre $b_\sigma$ y $w_j$.} Por el primer lema,
        $w_j\in\sigma$, luego, por el lema del baricentro,
        \[
          |w_j-b_\sigma|\leq\frac{n}{n+1}\diam\sigma.
        \]
      \item \alert{Entre $w_j$ y $w_k$.} Por hipótesis de inducción
        aplicada a $\partial_i\sigma$, que tiene dimensión $n-1$,
        \[
          |w_j-w_k|\leq\diam[w_0,\ldots,w_{n-1}]
          \leq\frac{n-1}{n}\diam \partial_i\sigma
          \leq\frac{n}{n+1}\diam\sigma,
        \]
        pues $\diam \partial_i\sigma\leq\diam\sigma$ y
        $\frac{n-1}{n}\leq\frac{n}{n+1}$, ya que $(n-1)(n+1)=n^2-1<n^2$.
    \end{itemize}

    \medskip

    Así, todas las distancias entre vértices de $\tau$ son
    $\leq\frac{n}{n+1}\diam\sigma$, y por tanto también lo es
    $\diam\tau$.
  \end{proof}
\end{frame}

\begin{frame}
  \frametitle{Subdivisión iterada y diámetro}

  \begin{corollary}
    Sea $\sigma$ un símplice generalizado de dimensión $n$. Para todo
    $k\geq0$ y todo $\tau\in\sd^k(\sigma)$,
    \[
      \diam\tau\leq\Bigl(\frac{n}{n+1}\Bigr)^k\diam\sigma.
    \]
    En particular, el máximo de los diámetros de los elementos de
    $\sd^k(\sigma)$ tiende a $0$ cuando $k\to\infty$.
  \end{corollary}

  \begin{proof}
    Por inducción en $k$, aplicando la proposición a cada elemento de
    $\sd^{k-1}(\sigma)$, que sigue teniendo dimensión $n$.
  \end{proof}
  
\end{frame}

\section{Cadenas singulares lineales}

% ------------------------------------------------------------------------
\subsection{Símplices singulares lineales}

\begin{frame}
  \frametitle{Símplices singulares lineales}

  En toda esta parte, $Y\subseteq\R^N$ es un subconjunto \alert{convexo}.

  \begin{definition}
    Un \alert{símplice singular lineal} de dimensión $n$ en $Y$ es una
    aplicación $\sigma\colon\simplex{n}\to Y$ que es la restricción de una
    aplicación lineal $\R^{n+1}\to\R^N$.
  \end{definition}

  Por linealidad, $\sigma$ está determinado por los vértices
  $v_i=\sigma(e_i)\in Y$, y
  \[
    \sigma(t_0,\ldots,t_n)=t_0v_0+\cdots+t_nv_n.
  \]
  Recíprocamente, como $Y$ es convexo, cualesquiera $v_0,\ldots,v_n\in Y$
  definen así un símplice singular lineal. Los símplices singulares lineales
  en $Y$ son \alert{lo mismo} que las listas de $n+1$ puntos de $Y$, y su
  imagen es el símplice generalizado $[v_0,\ldots,v_n]\subseteq Y$. Por eso
  escribimos
  \[
    \sigma=[v_0,\ldots,v_n].
  \]
\end{frame}

\begin{frame}
  \frametitle{Ejemplos}

  \begin{itemize}
    \item La identidad de $\simplex{n}$ es el símplice singular lineal
      \[
        \id_{\simplex{n}}=[e_0,\ldots,e_n]\colon
        \simplex{n}\longrightarrow\simplex{n}.
      \]
    \item Las inclusiones de las caras
      $\delta_i\colon\simplex{n-1}\to\simplex{n}$ son
      $\delta_i=[e_0,\ldots,\widehat{e}_i,\ldots,e_n]$.
    \item Si $\sigma=[v_0,\ldots,v_n]$, la cara $\sigma\delta_i$ es el
      símplice singular lineal $\partial_i\sigma=[v_0,\ldots,\widehat{v}_i,\ldots,v_n]$.
  \end{itemize}
\end{frame}

% ------------------------------------------------------------------------
\subsection{El complejo de cadenas lineales}

\begin{frame}
  \frametitle{Cadenas singulares lineales}

  \begin{definition}
    El grupo de \alert{cadenas singulares lineales}
    $\LC_n(Y)\subseteq\Sing_n(Y)$ es el subgrupo generado por los
    símplices singulares lineales de dimensión $n$ en $Y$.
  \end{definition}

  \medskip

  Es abeliano libre, con base los símplices singulares lineales en $Y$.
  Las caras de un símplice singular lineal son símplices singulares
  lineales, luego el borde de una cadena singular lineal es una cadena singular
  lineal: $\LC_\bullet(Y)$ es un subcomplejo de $\Sing_\bullet(Y)$, el
  \alert{complejo de cadenas singulares lineales} de $Y$.
\end{frame}

% ------------------------------------------------------------------------
\subsection{La versión aumentada}

\begin{frame}
  \frametitle{Cadenas singulares lineales aumentadas}

  \begin{definition}
    El complejo \alert{aumentado} $\widetilde{\LC}_\bullet(Y)$ coincide con
    $\LC_\bullet(Y)$ en dimensiones $n\geq0$, y
    $\widetilde{\LC}_{-1}(Y)=\Z\,[\;]$, con
    $\partial[v_0]=[\;]$ para todo $v_0\in Y$.
  \end{definition}

  \medskip

  Es un subcomplejo del complejo de cadenas singulares aumentado
  $\widetilde{\Sing}_\bullet(Y)$. Convenimos que $\partial=0$ en grado $-1$.
\end{frame}

\section{Subdivisión de cadenas lineales}

\begin{frame}
  \frametitle{Subdivisión de cadenas lineales}

  % Pendiente: S en LC(Y) definido recursivamente por S(σ)=b_σ·S(∂σ), morfismo de complejos y homotopía con la identidad.
\end{frame}

\section{El operador subdivisión}

\begin{frame}
  \frametitle{El operador subdivisión}

  % Pendiente: S en cadenas lineales y singulares, morfismo de complejos, homotopía de cadenas con la identidad.
\end{frame}

\section{Iteración de la subdivisión}

% ------------------------------------------------------------------------
\subsection{Iteración de homotopías}

\begin{frame}[allowframebreaks]
  \frametitle{Potencias de un morfismo homotópico a la identidad}

  \begin{lemma}
    Sea $C_\bullet$ un complejo de cadenas, $\Sd\colon C_\bullet\to C_\bullet$ un morfismo de complejos homotópico a la identidad y $T$ una
    homotopía de cadenas concreta, $\partial T+T\partial=\id-\Sd$. Para cada
    $k\geq0$, sea
    \[
      T_k=\sum_{j=0}^{k-1}T\Sd^j\colon C_n\to C_{n+1}.
    \]
    En particular, $T_0=0$ por ser una suma vacía.
    Entonces $\partial T_k+T_k\partial=\id-\Sd^k$. En particular, $\Sd^k$
    es homotópico a la identidad.
  \end{lemma}

  El subíndice de $T_k$ indica el número de iteraciones, no el grado.
  Directamente de la definición, si $m'\leq m$,
  \[
    T_m-T_{m'}=\sum_{j=m'}^{m-1}T\Sd^j.
  \]

  \framebreak

  \begin{proof}
    Como $\Sd$ es un morfismo de complejos, $\Sd^j\partial=\partial\Sd^j$
    para todo $j\geq0$. Por tanto, para cada $j$,
    \[
      \partial T\Sd^j+T\Sd^j\partial
      =\partial T\Sd^j+T\partial\Sd^j
      =(\partial T+T\partial)\Sd^j
      =(\id-\Sd)\Sd^j
      =\Sd^j-\Sd^{j+1}.
    \]
    Sumando para $j=0,\ldots,k-1$, vemos que en el miembro derecho se producen multitud de cancelaciones entre sumandos, excepto del primer y último término, y obtenemos:
    \[
      \partial T_k+T_k\partial
      =\sum_{j=0}^{k-1}\bigl(\Sd^j-\Sd^{j+1}\bigr)
      =\id-\Sd^k.
    \]
    Si $k=0$, ambos miembros son nulos, pues $T_0=0$ y $\Sd^0=\id$.
  \end{proof}
\end{frame}

\section{Cadenas subordinadas a un recubrimiento}

\begin{frame}
  \frametitle{Cadenas subordinadas a un recubrimiento}

  % Pendiente: complejo C^U(X) de símplices contenidos en algún U_i, con X = unión de los interiores.
\end{frame}

\section{El teorema de las cadenas pequeñas}

% ------------------------------------------------------------------------
\subsection{Enunciado}

\begin{frame}
  \frametitle{El teorema de las cadenas pequeñas}

  \begin{theorem}
    Sea $\cU$ un recubrimiento de un espacio $X$ e
    $\iota\colon\CU_\bullet(X)\hookrightarrow\Sing_\bullet(X)$ la inclusión.
    Existen un morfismo de complejos
    $\rho\colon\Sing_\bullet(X)\to\CU_\bullet(X)$ y homomorfismos
    $D\colon\Sing_n(X)\to\Sing_{n+1}(X)$ tales que
    \begin{enumerate}
      \item $\rho\iota=\id$, es decir, $\rho$ fija $\CU_\bullet(X)$;
      \item $\partial D+D\partial=\id-\iota\rho$;
      \item $D$ se anula sobre $\CU_\bullet(X)$.
    \end{enumerate}
  \end{theorem}

  En particular, $\iota$ es una equivalencia de homotopía de cadenas con
  inversa homotópica~$\rho$.

  \medskip

  La demostración se hará en varias etapas, usando la subdivisión iterada y el lema del número de Lebesgue.
\end{frame}


% ------------------------------------------------------------------------
\subsection{El número de subdivisiones}

\begin{frame}[allowframebreaks]
  \frametitle{El número de subdivisiones de un símplice}

  Sea $\sigma\colon\simplex{n}\to X$ un símplice singular. Por la
  proposición de subdivisión suficiente, existe un menor entero $k(\sigma)\geq0$ tal que
  $\Sd^i\sigma\in\CU_n(X)$ para todo $i\geq k(\sigma)$.

  \begin{definition}
    Definimos $m(\sigma)\geq0$ por recurrencia en $n$:
    \[
      m(\sigma)=
      \begin{cases}
        k(\sigma), & n=0,\\
        \max\bigl\{k(\sigma),m(\partial_0\sigma),\ldots,m(\partial_n\sigma)\bigr\},
          & n\geq1.
      \end{cases}
    \]
  \end{definition}

  \framebreak

  \begin{lemma}
    Para todo símplice singular $\sigma\colon\simplex{n}\to X$:
    \begin{enumerate}
      \item $\Sd^i\sigma\in\CU_n(X)$ para todo $i\geq m(\sigma)$;
      \item si $n\geq1$, $m(\partial_j\sigma)\leq m(\sigma)$ para todo $j$;
      \item si $\sigma\in\CU_n(X)$, entonces $m(\sigma)=0$.
    \end{enumerate}
  \end{lemma}

  \begin{proof}
    Las dos primeras propiedades se siguen de la definición, ya que
    $m(\sigma)\geq k(\sigma)$ y $m(\sigma)\geq m(\partial_j\sigma)$.

    \medskip

    La tercera, por inducción en $n$. Como $\Sd$ preserva
    $\CU_\bullet(X)$, $\Sd^i\sigma\in\CU_n(X)$ para todo $i\geq0$, luego
    $k(\sigma)=0$. Si $n\geq1$, las caras $\partial_j\sigma$ también están en
    $\CU_{n-1}(X)$, y por hipótesis de inducción $m(\partial_j\sigma)=0$.
    Así, $m(\sigma)=0$.
  \end{proof}
\end{frame}

% ------------------------------------------------------------------------
\subsection{Construcción de la retracción}

\begin{frame}[allowframebreaks]
  \frametitle{Definición de la homotopía y de la retracción}

  \begin{definition}
    Sea $D\colon\Sing_n(X)\to\Sing_{n+1}(X)$ el homomorfismo dado sobre los
    símplices singulares por
    \[
      D\sigma=T_{m(\sigma)}\sigma,
    \]
    y sea \[\rho=\id-\partial D-D\partial\colon\Sing_n(X)\to\Sing_n(X).\]
  \end{definition}

  \framebreak

  \begin{lemma}
    El homomorfismo $\rho\colon\Sing_\bullet(X)\to\Sing_\bullet(X)$ es un morfismo de complejos, $\partial\rho=\rho\partial$.
  \end{lemma}

  \begin{proof}
    Como $\partial^2=0$,
    \[
      \partial\rho
      = \partial(\id-\partial D-D\partial)
      =\partial-\partial D\partial
      \qquad\text{y}\qquad
      \rho\partial
      =(\id-\partial D-D\partial)\partial
      =\partial-\partial D\partial.
    \]
  \end{proof}

  Por definición, $\partial D+D\partial=\id-\rho$. Falta ver que $\rho$
  toma valores en $\CU_\bullet(X)$ y las demás propiedades.
\end{frame}

\begin{frame}
  \frametitle{La homotopía se anula sobre las cadenas subordinadas}

  \begin{lemma}
    El homomorfismo $D$ se anula sobre $\CU_\bullet(X)$ y $\rho$ fija
    $\CU_\bullet(X)$.
  \end{lemma}

  \begin{proof}
    Sea $\sigma\colon\simplex{n}\to X$ un símplice singular subordinado a $\cU$. Por un lema anterior,
    $m(\sigma)=0$, luego $D\sigma=T_0\sigma=0$. Esto prueba que $D$ se anula
    sobre $\CU_\bullet(X)$.

    \medskip

    Además, $\partial\sigma\in\CU_{n-1}(X)$, así que también
    $D\partial\sigma=0$, y por tanto
    \[
      \rho\sigma=\sigma-\partial D\sigma-D\partial\sigma=\sigma.
    \]
  \end{proof}
\end{frame}

\begin{frame}[allowframebreaks]
  \frametitle{La retracción toma valores en las cadenas subordinadas}

  \begin{lemma}
    Se tiene que $\rho\bigl(\Sing_n(X)\bigr)\subseteq\CU_n(X)$.
  \end{lemma}

  \begin{proof}
    Sea $\sigma\colon\simplex{n}\to X$ un símplice singular. Si $n=0$, la
    imagen de $\sigma$ es un punto, que está en algún $U_i$, luego
    $\sigma\in\CU_0(X)$ y $\rho\sigma=\sigma$ por el lema anterior.

    \medskip

    Si $n\geq1$, sean $m=m(\sigma)$ y $m_j=m(\partial_j\sigma)$, que cumplen
    $m_j\leq m$. Por el lema de iteración de homotopías,
    \[
      \partial D\sigma=\partial T_m\sigma
      =\sigma-\Sd^m\sigma-T_m\partial\sigma.
    \]

    \framebreak

    Sustituyendo en la definición de $\rho$ y usando que
    $\partial\sigma=\sum_j(-1)^j\partial_j\sigma$,
    \begin{align*}
      \rho\sigma
      &=\sigma-\partial D\sigma-D\partial\sigma
       =\Sd^m\sigma+T_m\partial\sigma-D\partial\sigma\\
      &=\Sd^m\sigma+\sum_{j=0}^n(-1)^j\bigl(T_m-T_{m_j}\bigr)(\partial_j\sigma)\\
      &=\Sd^m\sigma+\sum_{j=0}^n(-1)^j\sum_{i=m_j}^{m-1}T\Sd^i(\partial_j\sigma).
    \end{align*}

    El primer término está en $\CU_n(X)$ porque $m=m(\sigma)$. En los
    demás, $i\geq m_j=m(\partial_j\sigma)$, luego
    $\Sd^i(\partial_j\sigma)\in\CU_{n-1}(X)$, y como $T$ preserva las cadenas
    subordinadas, $T\Sd^i(\partial_j\sigma)\in\CU_n(X)$. Por tanto
    $\rho\sigma\in\CU_n(X)$.
  \end{proof}
\end{frame}

% ------------------------------------------------------------------------
\subsection{Consecuencias}

\begin{frame}
  \frametitle{Demostración del teorema}

  Los lemas anteriores prueban el teorema:
  \begin{itemize}
    \item $\rho\colon \Sing_\bullet(X)\to\CU_\bullet(X)$ es un morfismo de complejos;
    \item $\rho$ fija $\CU_\bullet(X)$, es decir, $\rho\iota=\id$;
    \item $\partial D+D\partial=\id-\iota\rho$, por definición de $\rho$;
    \item $D$ se anula sobre $\CU_\bullet(X)$.
  \end{itemize}
\end{frame}

\begin{frame}
  \frametitle{Consecuencia en homología}

  \begin{corollary}
    La inclusión $\iota\colon\CU_\bullet(X)\hookrightarrow\Sing_\bullet(X)$
    induce isomorfismos
    \[
      H_n\bigl(\CU_\bullet(X)\bigr)\cong H_n(X)
    \]
    para todo $n$, con inverso inducido por $\rho$.
  \end{corollary}

  \medskip

  Es decir, la homología singular puede calcularse usando solo símplices
  subordinados a un recubrimiento cualquiera.
\end{frame}

\begin{frame}[allowframebreaks]
  \frametitle{Versión reducida}

  El complejo aumentado $\tCU_\bullet(X)$ coincide con $\CU_\bullet(X)$
  en grados $n\geq0$ y con $\widetilde{\Sing}_{-1}(X)=\Z\,[\;]$ en grado
  $-1$. Es un subcomplejo de $\widetilde{\Sing}_\bullet(X)$.

  \begin{corollary}
    La inclusión $\tCU_\bullet(X)\hookrightarrow\widetilde{\Sing}_\bullet(X)$
    induce isomorfismos
    \[
      H_n\bigl(\tCU_\bullet(X)\bigr)\cong\widetilde{H}_n(X)
    \]
    para todo $n$.
  \end{corollary}

  \framebreak

  \begin{proof}
    Extendemos $\rho$ por la identidad y $D$ por $0$ en grado $-1$. Todo
    $0$-símplice tiene imagen en un punto, que está en algún $U_i$, así que
    $\CU_0(X)=\Sing_0(X)$ y $\rho$ es la identidad en grado $0$. Por tanto
    $\rho$ conmuta con el borde $\Sing_0(X)\to\widetilde{\Sing}_{-1}(X)$, y
    es un morfismo de complejos aumentados.

    \medskip

    Las identidades $\rho\iota=\id$ y $\partial D+D\partial=\id-\iota\rho$
    se cumplen en grados $n\geq1$ como antes. En grado $0$ también, pues
    $D\partial\sigma=D[\;]=0$, y en grado $-1$ ambos miembros son nulos.
    Así, la inclusión es una equivalencia de homotopía de cadenas.
  \end{proof}
\end{frame}

\section{El teorema de excisión}

% ------------------------------------------------------------------------
\subsection{Enunciado}

\begin{frame}
  \frametitle{El teorema de excisión}

  Para $A\subseteq X$, las cadenas relativas son
  $\Sing_\bullet(X,A)=\Sing_\bullet(X)/\Sing_\bullet(A)$, y la homología
  relativa es $H_n(X,A)=H_n\bigl(\Sing_\bullet(X,A)\bigr)$.

  \begin{theorem}[Excisión]
    Sean $Z\subseteq A\subseteq X$ con $\overline{Z}\subseteq\Int A$.
    Entonces la inclusión de pares
    $(X\setminus Z,A\setminus Z)\hookrightarrow(X,A)$ induce isomorfismos
    \[
      H_n(X\setminus Z,A\setminus Z)\cong H_n(X,A)
    \]
    para todo $n$.
  \end{theorem}

  Es decir, podemos \alert{suprimir} $Z$ sin alterar la homología
  relativa, siempre que $Z$ esté bien dentro de $A$.
\end{frame}

\begin{frame}[allowframebreaks]
  \frametitle{Una formulación equivalente}

  \begin{theorem}[Excisión, segunda forma]
    Sean $A,B\subseteq X$ con $\Int A\cup\Int B=X$. Entonces la inclusión
    de pares $(B,A\cap B)\hookrightarrow(X,A)$ induce isomorfismos
    \[
      H_n(B,A\cap B)\cong H_n(X,A)
    \]
    para todo $n$.
  \end{theorem}

  \framebreak

  \begin{proposition}
    La segunda forma implica la primera.
  \end{proposition}

  \begin{proof}
    Sean $Z\subseteq A\subseteq X$ con $\overline{Z}\subseteq\Int A$, y sea
    $B=X\setminus Z$. Entonces
    \[
      \Int B=X\setminus\overline{Z}\supseteq X\setminus\Int A,
    \]
    luego $\Int A\cup\Int B=X$. Además, $A\cap B=A\setminus Z$, así que la
    inclusión $(B,A\cap B)\hookrightarrow(X,A)$ es la del enunciado de la
    primera forma.
  \end{proof}

  \medskip

  Demostraremos, por tanto, la segunda forma. La hipótesis dice
  exactamente que $\cU=\{A,B\}$ es un recubrimiento de $X$.
\end{frame}

% ------------------------------------------------------------------------
\subsection{Demostración}

\begin{frame}
  \frametitle{Primer paso: un isomorfismo de complejos}

  Sea $\cU=\{A,B\}$, de modo que
  $\CU_\bullet(X)=\Sing_\bullet(A)+\Sing_\bullet(B)$.

  \begin{lemma}
    La inclusión induce un isomorfismo de complejos
    \[
      \Sing_\bullet(B)/\Sing_\bullet(A\cap B)
      \xrightarrow{\ \cong\ }
      \CU_\bullet(X)/\Sing_\bullet(A).
    \]
  \end{lemma}

  \begin{proof}
    Es el segundo teorema de isomorfía,
    $B'/(A'\cap B')\cong(A'+B')/A'$, aplicado en cada grado a
    $A'=\Sing_n(A)$ y $B'=\Sing_n(B)$, teniendo en cuenta que
    $\Sing_n(A)\cap\Sing_n(B)=\Sing_n(A\cap B)$. El isomorfismo es inducido
    por la inclusión, luego conmuta con el borde.
  \end{proof}
\end{frame}

\begin{frame}[allowframebreaks]
  \frametitle{Segundo paso: cadenas pequeñas relativas}

  \begin{lemma}
    La inclusión induce isomorfismos en homología
    \[
      H_n\bigl(\CU_\bullet(X)/\Sing_\bullet(A)\bigr)
      \xrightarrow{\ \cong\ }
      H_n\bigl(\Sing_\bullet(X)/\Sing_\bullet(A)\bigr).
    \]
  \end{lemma}

  \begin{proof}
    Sean $\rho$ y $D$ los del teorema de las cadenas pequeñas. Como
    $\Sing_\bullet(A)\subseteq\CU_\bullet(X)$, se tiene que $\rho$ fija
    $\Sing_\bullet(A)$ y $D$ se anula sobre $\Sing_\bullet(A)$. Por tanto
    $\iota$, $\rho$ y $D$ pasan al cociente por $\Sing_\bullet(A)$:
    \[
      \bar\iota\colon\CU_\bullet(X)/\Sing_\bullet(A)
        \rightleftarrows\Sing_\bullet(X)/\Sing_\bullet(A)\colon\bar\rho,
      \qquad
      \bar D\colon\Sing_n(X,A)\to\Sing_{n+1}(X,A).
    \]

    \framebreak

    Las identidades del teorema de las cadenas pequeñas pasan al cociente:
    \[
      \bar\rho\bar\iota=\id,
      \qquad
      \partial\bar D+\bar D\partial=\id-\bar\iota\bar\rho.
    \]
    Así, $\bar\iota$ es una equivalencia de homotopía de cadenas con inversa
    $\bar\rho$, y por tanto induce isomorfismos en homología.
  \end{proof}
\end{frame}

\begin{frame}
  \frametitle{Conclusión}

  \begin{proof}[Demostración de la excisión]
    La inclusión $\Sing_\bullet(B,A\cap B)\to\Sing_\bullet(X,A)$ inducida por
    $(B,A\cap B)\hookrightarrow(X,A)$ es la composición
    \[
      \Sing_\bullet(B)/\Sing_\bullet(A\cap B)
      \xrightarrow{\ \cong\ }
      \CU_\bullet(X)/\Sing_\bullet(A)
      \xrightarrow{\ \bar\iota\ }
      \Sing_\bullet(X)/\Sing_\bullet(A).
    \]
    El primer morfismo es un isomorfismo de complejos por el primer paso, y
    el segundo induce isomorfismos en homología por el segundo paso. Por
    tanto, la composición induce isomorfismos
    $H_n(B,A\cap B)\cong H_n(X,A)$.
  \end{proof}
\end{frame}

% ------------------------------------------------------------------------
\subsection{Versión reducida}

\begin{frame}
  \frametitle{Excisión en homología reducida}

  Para $A\subseteq X$, la homología relativa reducida
  $\widetilde{H}_n(X,A)$ es la homología de
  $\widetilde{\Sing}_\bullet(X)/\widetilde{\Sing}_\bullet(A)$.

  \begin{corollary}
    En las hipótesis de la excisión, si además $A\setminus Z\neq\varnothing$,
    la inclusión induce isomorfismos
    \[
      \widetilde{H}_n(X\setminus Z,A\setminus Z)\cong\widetilde{H}_n(X,A)
    \]
    para todo $n$.
  \end{corollary}

  \begin{proof}
    Si $A\neq\varnothing$, los complejos aumentados de $X$ y de $A$ coinciden
    en grado $-1$, luego
    $\widetilde{\Sing}_\bullet(X)/\widetilde{\Sing}_\bullet(A)
    =\Sing_\bullet(X)/\Sing_\bullet(A)$ y
    $\widetilde{H}_n(X,A)=H_n(X,A)$. Lo mismo vale para el par
    $(X\setminus Z,A\setminus Z)$, pues $A\setminus Z\neq\varnothing$. Basta
    aplicar el teorema de excisión.
  \end{proof}
\end{frame}

\section{La sucesión de Mayer--Vietoris}

% ------------------------------------------------------------------------
\subsection{Enunciado}

\begin{frame}
  \frametitle{La sucesión de Mayer--Vietoris}

  Sean $A,B\subseteq X$ con $\Int A\cup\Int B=X$, y sean
  \[
    i_A\colon A\cap B\hookrightarrow A,\quad
    i_B\colon A\cap B\hookrightarrow B,\quad
    j_A\colon A\hookrightarrow X,\quad
    j_B\colon B\hookrightarrow X
  \]
  las inclusiones.

  \begin{theorem}[Mayer--Vietoris]
    Existe una sucesión exacta larga
    \begin{multline*}
      \cdots\to H_n(A\cap B)\xrightarrow{\left(\begin{smallmatrix}{i_A}_*\\-{i_B}_*\end{smallmatrix}\right)}H_n(A)\oplus H_n(B)
      \xrightarrow{\left(\begin{smallmatrix}{j_A}_* & {j_B}_*\end{smallmatrix}\right)}H_n(X)\\
      \xrightarrow{\ \delta\ }H_{n-1}(A\cap B)\to\cdots\to H_0(X)\to0.
    \end{multline*}
  \end{theorem}

  Este teorema es consecuencia del lema siguiente y del teorema de las
  cadenas pequeñas.
\end{frame}

% ------------------------------------------------------------------------
\subsection{Demostración}

\begin{frame}[allowframebreaks]
  \frametitle{Una sucesión exacta corta de complejos}

  Sea $\cU=\{A,B\}$, de modo que
  $\CU_\bullet(X)=\Sing_\bullet(A)+\Sing_\bullet(B)$.

  \begin{lemma}
    La sucesión de complejos
    \[
      0\to\Sing_\bullet(A\cap B)\xrightarrow{\left(\begin{smallmatrix}{i_A}_\#\\-{i_B}_\#\end{smallmatrix}\right)}
      \Sing_\bullet(A)\oplus\Sing_\bullet(B)\xrightarrow{\left(\begin{smallmatrix}{j_A}_\# & {j_B}_\#\end{smallmatrix}\right)}
      \CU_\bullet(X)\to0
    \]
    es exacta.
  \end{lemma}

  \begin{proof}
    Obvio porque $\CU_\bullet(X)=\Sing_\bullet(A)+\Sing_\bullet(B)$ y $\Sing_\bullet(A)\cap\Sing_\bullet(B)=\Sing_\bullet(A\cap B)$.
  \end{proof}
\end{frame}

\begin{frame}
  \frametitle{El morfismo de conexión}

  \begin{proposition}
    Sea $\alpha\in H_n(X)$. Existe un ciclo $x+y$ que lo representa con
    $x\in\Sing_n(A)$ e $y\in\Sing_n(B)$, y para cualquiera de ellos
    \[
      \delta\alpha=[\partial x]\in H_{n-1}(A\cap B).
    \]
  \end{proposition}

  \begin{proof}
    Por el teorema de las cadenas pequeñas, $\alpha$ tiene un representante
    en $\CU_n(X)=\Sing_n(A)+\Sing_n(B)$. Por la construcción del morfismo de
    conexión, se eleva $x+y\in \CU_n(X)$ a $(x,y)\in \Sing_n(A)\oplus\Sing_n(B)$, se toma su borde $(\partial
    x,\partial y)$ y se busca su preimagen por $\left(\begin{smallmatrix}{i_A}_\#\\-{i_B}_\#\end{smallmatrix}\right)$. Como
    $\partial x+\partial y=\partial(x+y)=0$, la cadena
    $\partial x=-\partial y$ está en $\Sing_{n-1}(A\cap B)$ y su imagen es
    $(\partial x,-\partial x)=(\partial x,\partial y)$.
  \end{proof}
\end{frame}

% ------------------------------------------------------------------------
\subsection{Versión reducida}

\begin{frame}[allowframebreaks]
  \frametitle{Mayer--Vietoris en homología reducida}

  \begin{theorem}
    Si además $A\cap B\neq\varnothing$, hay una sucesión exacta larga
    \begin{multline*}
      \cdots\to\widetilde{H}_n(A\cap B)\xrightarrow{\left(\begin{smallmatrix}{i_A}_*\\-{i_B}_*\end{smallmatrix}\right)}
      \widetilde{H}_n(A)\oplus\widetilde{H}_n(B)
      \xrightarrow{\left(\begin{smallmatrix}{j_A}_* & {j_B}_*\end{smallmatrix}\right)}\widetilde{H}_n(X)\\
      \xrightarrow{\ \delta\ }\widetilde{H}_{n-1}(A\cap B)\to\cdots
      \to\widetilde{H}_0(X)\to0,
    \end{multline*}
    con el morfismo de conexión $\delta$ dado por la misma fórmula.
  \end{theorem}

  \framebreak

  \begin{proof}
    Sea $\tCU_\bullet(X)$ el complejo aumentado
    de las cadenas subordinadas. La sucesión del lema se extiende a los
    complejos aumentados:
    \[
      0\to\widetilde{\Sing}_\bullet(A\cap B)\xrightarrow{\left(\begin{smallmatrix}{i_A}_\#\\-{i_B}_\#\end{smallmatrix}\right)}
      \widetilde{\Sing}_\bullet(A)\oplus\widetilde{\Sing}_\bullet(B)
      \xrightarrow{\left(\begin{smallmatrix}{j_A}_\# & {j_B}_\#\end{smallmatrix}\right)}\tCU_\bullet(X)\to0.
    \]
    En grado $-1$ es $0\to\Z\to\Z\oplus\Z\to\Z\to0$, con $x\mapsto(x,-x)$ y
    $(y,z)\mapsto y+z$, que es exacta.

    \medskip

    Su sucesión exacta larga, junto con la versión reducida del teorema de
    las cadenas pequeñas, da la sucesión del enunciado. Termina en
    $\widetilde{H}_0(X)\to0$ porque
    $\widetilde{H}_{-1}(A\cap B)=0$, ya que $A\cap B\neq\varnothing$.
  \end{proof}
\end{frame}

\section{Versiones para $CW$-complejos}

% ------------------------------------------------------------------------
\subsection{Subcomplejos}

\begin{frame}[allowframebreaks]
  \frametitle{Subespacios con buenos entornos}

  En esta sección, $A,B\subseteq X$ son subespacios con $X=A\cup B$ para los
  que existen abiertos $U\supseteq A$ y $V\supseteq B$ tales que
  \[
    A\subseteq U,\qquad B\subseteq V,\qquad A\cap B\subseteq U\cap V
  \]
  son retractos por deformación. Nótese que $X\supseteq U\cup V\supseteq A\cup B=X$,
  luego $\{U,V\}$ es un recubrimiento de $X$.

  \framebreak

  El caso principal es el de los $CW$-complejos.

  \begin{proposition}[Hatcher, Prop.~A.5]
    Sea $X$ un $CW$-complejo. A cada subcomplejo $A\subseteq X$ se le puede
    asociar un entorno abierto $N(A)\subseteq X$ de modo que
    \begin{enumerate}
      \item $A$ es un retracto por deformación de $N(A)$;
      \item $N(A)\cap N(B)=N(A\cap B)$ para cualesquiera subcomplejos $A,B$.
    \end{enumerate}
  \end{proposition}

  Por tanto, si $X$ es un $CW$-complejo y
  $A,B\subseteq X$ son subcomplejos con $X=A\cup B$, los abiertos $U=N(A)$
  y $V=N(B)$ cumplen la hipótesis, ya que $U\cap V=N(A\cap B)$.
\end{frame}

\begin{frame}[allowframebreaks]
  \frametitle{Reemplazar $A$ y $B$ por $U$ y $V$}

  Como los retractos de deformación son equivalencias homotópicas, el siguiente resultado es inmediato.

  \begin{lemma}
    Sean $A,B\subseteq X$ con $X=A\cup B$, y $U\supseteq A$, $V\supseteq B$
    abiertos como antes. Las inclusiones inducen isomorfismos
    \begin{gather*}
      H_n(A)\cong H_n(U),\qquad H_n(B)\cong H_n(V),\qquad
      H_n(A\cap B)\cong H_n(U\cap V),\\
      H_n(X,A)\cong H_n(X,U),\qquad H_n(B,A\cap B)\cong H_n(V,U\cap V),
    \end{gather*}
    y lo mismo en homología reducida.
  \end{lemma}

\end{frame}

% ------------------------------------------------------------------------
\subsection{Escisión y Mayer--Vietoris}

\begin{frame}[allowframebreaks]
  \frametitle{Escisión}

  \begin{theorem}
    Sean $A,B\subseteq X$ como antes; por ejemplo, subcomplejos de un
    $CW$-complejo con $X=A\cup B$. Entonces la inclusión de pares
    $(B,A\cap B)\hookrightarrow(X,A)$ induce isomorfismos
    \[
      H_n(B,A\cap B)\cong H_n(X,A)
    \]
    para todo $n$.
  \end{theorem}

  \framebreak

  \begin{proof}
    Las inclusiones de pares dan un diagrama conmutativo
    \[
      \begin{tikzcd}[ampersand replacement=\&]
        H_n(B,A\cap B)\arrow[r]\arrow{d}[swap]{\cong}
          \&H_n(X,A)\arrow{d}{\cong}\\
        H_n(V,U\cap V)\arrow{r}{\cong}\&H_n(X,U)
      \end{tikzcd}
    \]
    Las flechas verticales son isomorfismos por el lema anterior. La
    inferior lo es por el teorema de escisión aplicado al recubrimiento
    abierto $\{U,V\}$. Por tanto la
    superior también es un isomorfismo.
  \end{proof}
\end{frame}

\begin{frame}[allowframebreaks]
  \frametitle{Mayer--Vietoris}

  \begin{theorem}
    Sean $A,B\subseteq X$ como antes; por ejemplo, subcomplejos de un
    $CW$-complejo con $X=A\cup B$. Con las inclusiones $i_A,i_B,j_A,j_B$ de
    antes, existe una sucesión exacta larga
    \begin{multline*}
      \cdots\to H_n(A\cap B)\xrightarrow{\left(\begin{smallmatrix}{i_A}_*\\-{i_B}_*\end{smallmatrix}\right)}H_n(A)\oplus H_n(B)
      \xrightarrow{\left(\begin{smallmatrix}{j_A}_* & {j_B}_*\end{smallmatrix}\right)}H_n(X)\\
      \xrightarrow{\ \delta\ }H_{n-1}(A\cap B)\to\cdots\to H_0(X)\to0.
    \end{multline*}
    Si además $A\cap B\neq\varnothing$, lo mismo vale en homología reducida.
  \end{theorem}

  \framebreak

  \begin{proof}
    Como $\{U,V\}$ es un recubrimiento abierto de $X$, tenemos la sucesión
    de Mayer--Vietoris de $U$ y $V$. Por el
    lema anterior, las inclusiones inducen isomorfismos
    \[
      H_n(A\cap B)\cong H_n(U\cap V),\qquad
      H_n(A)\cong H_n(U),\qquad
      H_n(B)\cong H_n(V),
    \]
    también en homología reducida, 
    luego el teorema se sigue.
  \end{proof}
\end{frame}


\end{document}
